%%%%%%%%%%%%%%%%%%%%%%%%%%%%%%%%%%%%%%%%%%%%%%%%%%%%%%%%%%%%%
%% Examtask 2: buck-boost converter as smart voltage stabilizer for bicycle dynamo %%
%%%%%%%%%%%%%%%%%%%%%%%%%%%%%%%%%%%%%%%%%%%%%%%%%%%%%%%%%%%%%

\task{Buck-boost converter as smart voltage stabilizer for bicycle dynamo}

\taskGerman{Tief-Hochsetzsteller als intelligente Spannungsstabilisierung für Fahrraddynamos}

A buck-boost converter is to be designed for voltage stabilization in a bicycle dynamo.
The dynamo generator supplies a variable voltage. A control circuit adjusts the converter's duty cycle
so that the buck-boost converter provides a stable output voltage.

\begin{germanblock}
	Für ein Fahraddynamo soll ein Tief-Hochsetzsteller als Spannungsstabilisierung dimensioniert werden.
	Der Dynamogenerator liefert eine variable Spannung. Eine Reglungsschaltung steuert das Tastverhältnis des Konverters,
	so das der Tief-Hochsetzsteller eine stabile Ausgangsspannung liefert.
\end{germanblock}

\input{fig/Task02/figbuckboostConverter}

\begin{table}[ht]
	\centering  % Zentriert die Tabelle
	\begin{tabular}{llll}
		\toprule
		\multicolumn{2}{l}{\textbf{General parameters:}} & \multicolumn{2}{l}{\textbf{MOSFET:}}                                                                                                          \\
		Input voltage:                                   & $U_{\mathrm{1}} = \SI{3}{\volt} \ldots \SI{10}{\volt}$           & Transistor drain source voltage:    & $u_\mathrm{ds,on} = \SI{0.5}{\volt}$ \\
		Output voltage:                                  & $U_{\mathrm{2}} = \SI{5}{\volt}$                                 &                                     &                                      \\
		Output current range:                            & $I_{\mathrm{2}} = \SI{50}{\milli\ampere} \ldots \SI{1}{\ampere}$ & \multicolumn{2}{l}{\textbf{Diode:}}                                        \\
		Switching frequency:                             & $f_\mathrm{s} = \SI{100}{\kilo\hertz}$                           & Diode forward voltage:              & $u_\mathrm{fw} = \SI{0.8}{\volt}$    \\
		\midrule
		\multicolumn{4}{l}{\quad The diode forward voltage and transistor voltage loss must be taken into account.}                                                                                      \\
		\multicolumn{4}{l}{\quad The drain-source voltage is valid for both transistors in conduct case..}                                                                                               \\
		\multicolumn{4}{l}{\quad The output voltage can be considered as constant except of subtask 2.4.}                                                                                                \\
		\bottomrule
	\end{tabular}
	\caption{Parameters of the circuit.}  % Beschriftung der Tabelle
	\label{table:ex02_Parameters of the buck-boost converter}
\end{table}

\subtask{What are the minimum and maximum duty cycles that the control circuit has to request in case of $D = D_1 = D_2$?}{2}
\subtaskGerman{Welches minimale und maximale Tastverhältnis muss die Reglungsschaltung für den Fall $D = D_1 = D_2$ vorgeben.}

% Solution of subtask
\begin{solutionblock}
	The synchronized duty cycle for the buck-boost converter is calculated using the inductor current balance equation
	\begin{equation*}
		\Delta i_\mathrm{L,T_\mathrm{on}} \cdot T_\mathrm{on}=\Delta i_\mathrm{L,T_\mathrm{off}} \cdot T_\mathrm{off}.
	\end{equation*}
	If the transistors conduct $\Delta i_\mathrm{L,T_\mathrm{on}}$ results in:
	\begin{equation*}
		\Delta i_\mathrm{L,T_\mathrm{on}}=\frac{(U_\mathrm{1}-2u_\mathrm{ds,on}) \cdot T_\mathrm{on} }{L}.
	\end{equation*}
	If the transistors are blocking $\Delta i_\mathrm{L,T_\mathrm{off}}$ is calculated by:
	\begin{equation*}
		\Delta i_\mathrm{L,T_\mathrm{off}}=\frac{(U_\mathrm{2}+ 2u_\mathrm{fw}) \cdot T_\mathrm{off}}{L}.
	\end{equation*}

	By substituting the currents within the inductor current balance equation, we obtain
	\begin{equation*}
		\frac{(U_\mathrm{1}-2u_\mathrm{ds,on}) \cdot T_\mathrm{on}}{L}=\frac{(U_\mathrm{2}+2u_\mathrm{fw}) \cdot T_\mathrm{off}}{L}.
	\end{equation*}

	Replacing $T_\mathrm{on}$ and $T_\mathrm{off}$ by the switching period yields to
	\begin{equation*}
		\frac{(U_\mathrm{1}-2u_\mathrm{ds,on}) \cdot T_\mathrm{s} \cdot D}{L} =
		\frac{(U_\mathrm{2}+2u_\mathrm{fw}) \cdot (1-D) \cdot T_\mathrm{s}}{L}.
	\end{equation*}

	After simplifying by canceling $T_\mathrm{s}$ and $L$, we obtain
	\begin{equation*}
		(U_\mathrm{1}-2u_\mathrm{ds,on}) \cdot D =
		(U_\mathrm{2}+2u_\mathrm{fw}) \cdot (1-D).
	\end{equation*}
	Solving this equation with respect to $D$ yields
	\begin{equation*}
		D = \frac{U_\mathrm{2}+2u_\mathrm{fw}}{U_\mathrm{1}-2u_\mathrm{ds,on}+U_\mathrm{2}+2u_\mathrm{fw}}.
	\end{equation*}
	Entering the minimum and maximum input voltage leads yields
	\begin{align*}
		D_\mathrm{max} & = \frac{\SI{5}{\volt}+ 2 \cdot \SI{0.8}{\volt}}{\SI{3}{\volt} - 2 \cdot \SI{0.5}{\volt} + \SI{5}{\volt}+ 2 \cdot \SI{0.8}{\volt}}
		= 0.77                                                                                                                                              \\
		D_\mathrm{min} & = \frac{\SI{5}{\volt}+ 2 \cdot \SI{0.8}{\volt}}{\SI{10}{\volt} - 2 \cdot \SI{0.5}{\volt} + \SI{5}{\volt}+ 2 \cdot \SI{0.8}{\volt}}
		= 0.42.
	\end{align*}
\end{solutionblock}

\subtask{Calculate the efficiency of the converter at the operation points $U_\mathrm{1}=\SI{3}{\volt}$, $U_\mathrm{1}=\SI{7.5}{\volt}$, and $U_\mathrm{1}=\SI{10}{\volt}$
	at maximum output power. Consider only the losses due to the constant voltage drops of the semiconductors. }{3}
\subtaskGerman{Berechne den Wirkungsgrand des Konverters an den Arbeitspunkten $U_\mathrm{1}=\SI{3}{\volt}$, $U_\mathrm{1}=\SI{7,5}{\volt}$, und $U_\mathrm{1}=\SI{10}{\volt}$
	bei maximaler Ausgangsleistung. Betrachten Sie nur die Verluste durch die Spannungsabfälle über die Halbleiter.}

% Solution of subtask
\begin{solutionblock}
	The power loss are be separated in 2 parts: Power loss of the transistors while $T_\mathrm{on}$-phase and power loss of the diodes while $T_\mathrm{off}$-phase.
	The average current of the inductance is the same for both parts and yields
	\begin{equation*}
		\overline{i}_\mathrm{L} = \frac{I_\mathrm{2}}{1-D}.
	\end{equation*}
	The duty cycle of operation point $U_\mathrm{1}=\SI{7.5}{\volt}$ results in
	\begin{equation*}
		D = \frac{U_\mathrm{2}+2u_\mathrm{fw}}{U_\mathrm{1}-2u_\mathrm{ds,on}+U_\mathrm{2}+2u_\mathrm{fw}}
		= \frac{\SI{5}{\volt}+ 2 \cdot \SI{0.8}{\volt}}{\SI{7.5}{\volt} - 2 \cdot \SI{0.5}{\volt}+\SI{0.8}{\volt + \SI{5}{\volt}+ 2 \cdot \SI{0.8}{\volt}}}
		= 0.5.
	\end{equation*}

	The power loss is calculated by
	\begin{equation*}
		P_\mathrm{loss}= P_\mathrm{loss, on}+P_\mathrm{loss, off}=\overline{i}_\mathrm{L} (2u_\mathrm{ds,on} \cdot D  + 2u_\mathrm{fw} \cdot (1-D)).
	\end{equation*}
	For the operation points we obtain
	\begin{align*}
		P_\mathrm{loss,U_\mathrm{1}=\SI{3}{\volt}}= \frac{\SI{1}{\ampere}}{1-0.77} (2 \cdot \SI{0.5}{\volt} \cdot 0.77  + 2 \cdot \SI{0.5}{\volt} \cdot (1-0.77))
		= \SI{4.9}{\watt}, \\
		P_\mathrm{loss,U_\mathrm{1}=\SI{7.5}{\volt}}= \frac{\SI{1}{\ampere}}{1-0.5} (2 \cdot \SI{0.5}{\volt} \cdot 0.5  + 2 \cdot \SI{0.5}{\volt} \cdot (1-0.5))
		= \SI{2.6}{\watt}, \\
		P_\mathrm{loss,U_\mathrm{1}=\SI{10}{\volt}}= \frac{\SI{1}{\ampere}}{1-0.42} (2 \cdot \SI{0.5}{\volt} \cdot 0.42  + 2 \cdot \SI{0.5}{\volt} \cdot (1-0.42))
		= \SI{2.3}{\watt}. \\
	\end{align*}
	The effciency is calculated by
	\begin{equation*}
		\eta = \frac{P_\mathrm{2}}{P_\mathrm{2}+P_\mathrm{loss}} = \frac{U_\mathrm{2} I_\mathrm{2}} {U_\mathrm{2} I_\mathrm{2}+P_\mathrm{loss}}.
	\end{equation*}
	For the operation points we obtain
	\begin{align*}
		\eta_\mathrm{U_\mathrm{1}=\SI{3}{\volt}} = \frac{\SI{5}{\volt} \cdot \SI{1}{\ampere}} {\SI{5}{\volt} \cdot \SI{1}{\ampere} + \SI{4.9}{\watt}} = 0.51,   \\
		\eta_\mathrm{U_\mathrm{1}=\SI{7.5}{\volt}} = \frac{\SI{5}{\volt} \cdot \SI{1}{\ampere}} {\SI{5}{\volt} \cdot \SI{1}{\ampere} + \SI{2.6}{\watt}} = 0.66, \\
		\eta_\mathrm{U_\mathrm{1}=\SI{10}{\volt}} = \frac{\SI{5}{\volt} \cdot \SI{1}{\ampere}} {\SI{5}{\volt} \cdot \SI{1}{\ampere} + \SI{2.3}{\watt}} = 0.68.  \\
	\end{align*}
\end{solutionblock}

\subtask{To avoid operation in discontinuous conduction mode (DCM) within the defined operating range, the minimum inductance  $L$ must be calculated.
	Evaluate the general equation for boundary conduction mode (BCM) to indicate, at which operation point the discontinuous conduction mode starts?  $D = D_1 = D_2$ still applies.}{2}
\subtaskGerman{ Um den Lückbetrieb innerhalb des definierten Arbeitsbereichs zu vermeiden, soll die minimale Induktivität
	$L$ berechnet werden.\\
	Werten Sie die allgemeine Gleichung für den Lückgrenzbetrieb aus, um den Arbeitspunkt zu bestimmen, an dem der Übergang zum Lückbetrieb auftritt.
	Berechnen Sie für diesen Arbeitspunkt die minimal erforderliche Induktivität um den Lückbetrieb zu vermeiden. $D = D_1 = D_2$ ist noch gültig.}

% Solution of subtask
\begin{solutionblock}
	The ripple current is calculated by
	\begin{equation*}
		\Delta i_\mathrm{L}=\frac{U_\mathrm{L,off} \cdot T_\mathrm{off}}{L}
		=\frac{(U_\mathrm{2}+2u_\mathrm{fw}) \cdot (1-D)}{L \cdot f_\mathrm{s}}.
	\end{equation*}
	Remark: The duty cycle depends on input voltage $U_\mathrm{1}$. Therefore the fix output voltage $U_\mathrm{2}$ is utilized.
	At boundary conduction mode (BCM) the ripple current yields
	\begin{equation*}
		\Delta i_\mathrm{L} = \frac{2 I_\mathrm{2}}{1-D}.
	\end{equation*}
	This leads to
	\begin{equation*}
		\frac{(U_\mathrm{2}+2u_\mathrm{fw}) \cdot (1-D)}{L \cdot f_\mathrm{s}} = \frac{2 I_\mathrm{2}}{1-D}.
	\end{equation*}
	Solving the equation with respect to $L$ leads to
	\begin{equation*}
		L = \frac{(U_\mathrm{2}+2u_\mathrm{fw}) \cdot (1-D)^2}{2 I_\mathrm{2} \cdot f_\mathrm{s}}.
	\end{equation*}
	The term $(1-D)^2$ express, that at minimum duty cycle the highest inductance is needed, to prevent discontinuous conduction mode.
	Therefore $D_\mathrm{min}$ is used delivering
	\begin{equation*}
		L = \frac{(\SI{5}{\volt}+2 \cdot \SI{0.8}{\volt}) \cdot (1-0.42)^2}{2 \cdot \SI{50}{\milli\ampere} \cdot \SI{100}{\kilo\hertz}}
		= \SI{219}{\micro\henry}.
	\end{equation*}

\end{solutionblock}

\subtask{ Now the control circuit regulates the transistors independent. It means, for low input voltage the boost mode is used and for high input voltage the buck mode is uses.
	How are the duty cycles $D_\mathrm{1}$ of transistor $T_\mathrm{1}$ and $D_\mathrm{2}$ of transistor $T_\mathrm{2}$ for the operating points?}{3}
\subtaskGerman{Nun regelt die Steuerung die Transistoren unabhängig voneinander. Das bedeutet, bei niedriger Eingangsspannung wird der Hochsetzsteller-Modus verwendet,
	bei hoher Eingangsspannung der Tiefsetzsteller-Modus. Wie lauten die Tastverhältnisse $D_\mathrm{1}$ des Transistors $T_\mathrm{1}$  und $D_\mathrm{2}$
	des Transistors $T_\mathrm{2}$ für die jeweiligen Betriebspunkte?}

\begin{solutionblock}
	In boost mode the duty cycle $D_\mathrm{1}=1$ so that $T_\mathrm{1}$ always conducts.
	The boost mode is active if $U_\mathrm{1}-u_\mathrm{ds,on}$ is less equal $U_\mathrm{2}+2u_\mathrm{fw}$. The break-even point is calculated by
	\begin{equation*}
		U_\mathrm{1}=U_\mathrm{2}+ u_\mathrm{fw} + u_\mathrm{ds,on}=\SI{5}{\volt}+\SI{0.8}{\volt}+ \SI{0.5}{\volt}=\SI{6.3}{\volt}.
	\end{equation*}
	The duty cycle for the boost mode is calculated using the inductor current balance equation
	\begin{equation*}
		\Delta i_\mathrm{L,T_\mathrm{on}} \cdot T_\mathrm{on}=\Delta i_\mathrm{L,T_\mathrm{off}} \cdot T_\mathrm{off}.
	\end{equation*}
	If the transistors conduct $\Delta i_\mathrm{L,T_\mathrm{on}}$ results in:
	\begin{equation*}
		\Delta i_\mathrm{L,T_\mathrm{on}}=\frac{(U_\mathrm{1}-2u_\mathrm{ds,on}) \cdot T_\mathrm{on} }{L}.
	\end{equation*}
	If the transistor $T_\mathrm{2}$ is blocking $\Delta i_\mathrm{L,T_\mathrm{off}}$ is calculated by:
	\begin{equation*}
		\Delta i_\mathrm{L,T_\mathrm{off}}=\frac{(U_\mathrm{2}+ u_\mathrm{fw} + u_\mathrm{ds,on}-U_\mathrm{1} ) \cdot T_\mathrm{off}}{L}.
	\end{equation*}
	By substituting the currents within the inductor current balance equation we obtain
	\begin{equation*}
		\frac{(U_\mathrm{1}-2u_\mathrm{ds,on}) \cdot T_\mathrm{on} }{L}=\frac{(U_\mathrm{2}+ u_\mathrm{fw} + u_\mathrm{ds,on}-U_\mathrm{1} ) \cdot T_\mathrm{off}}{L}.
	\end{equation*}
	Replacing $T_\mathrm{on}$ and $T_\mathrm{off}$ by the switching period yields to
	\begin{equation*}
		\frac{(U_\mathrm{1}-2u_\mathrm{ds,on}) \cdot T_\mathrm{s} \cdot D_\mathrm{2}}{L} =
		\frac{(U_\mathrm{2}+ u_\mathrm{fw} + u_\mathrm{ds,on}-U_\mathrm{1}) \cdot (1-D_\mathrm{2}) \cdot T_\mathrm{s}}{L}.
	\end{equation*}
	After simplifying by canceling $T_\mathrm{s}$ and $L$ we obtain
	\begin{equation*}
		(U_\mathrm{1}-2u_\mathrm{ds,on}) \cdot D_\mathrm{2} =
		(U_\mathrm{2}+ u_\mathrm{fw} + u_\mathrm{ds,on}-U_\mathrm{1}) \cdot (1-D_\mathrm{2}).
	\end{equation*}
	Solving this equation with respect to $D_\mathrm{2}$ yields
	\begin{equation*}
		D_\mathrm{2} = \frac{U_\mathrm{2}+ u_\mathrm{fw} + u_\mathrm{ds,on}-U_\mathrm{1}}{U_\mathrm{2}+ u_\mathrm{fw} - u_\mathrm{ds,on}}.
	\end{equation*}
	For the operation point $U_\mathrm{1}=\SI{3}{\volt}$ we obtain
	\begin{equation*}
		D_\mathrm{2,U_\mathrm{1}=\SI{3}{\volt}}=\frac{\SI{5}{\volt}+ \SI{0.8}{\volt} + \SI{0.5}{\volt}-\SI{3}{\volt}}{\SI{5}{\volt} + \SI{0.8}{\volt} - \SI{0.5}{\volt}}
		= 0.62.
	\end{equation*}
	In buck mode the duty cycle $D_\mathrm{2}=0$ so that $T_\mathrm{2}$ always blocks.
	The duty cycle for the buck mode is calculated using the inductor current balance equation
	\begin{equation*}
		\Delta i_\mathrm{L,T_\mathrm{on}} \cdot T_\mathrm{on}=\Delta i_\mathrm{L,T_\mathrm{off}} \cdot T_\mathrm{off}.
	\end{equation*}
	If the transistor $T_\mathrm{1}$ conducts $\Delta i_\mathrm{L,T_\mathrm{on}}$ results in:
	\begin{equation*}
		\Delta i_\mathrm{L,T_\mathrm{on}}=\frac{(U_\mathrm{1}-u_\mathrm{ds,on} - u_\mathrm{fw} - U_\mathrm{2}) \cdot T_\mathrm{on} }{L}.
	\end{equation*}
	If the transistor $T_\mathrm{1}$ is blocking $\Delta i_\mathrm{L,T_\mathrm{off}}$ is calculated by:
	\begin{equation*}
		\Delta i_\mathrm{L,T_\mathrm{off}}=\frac{(U_\mathrm{2} + 2 u_\mathrm{fw}) \cdot T_\mathrm{off}}{L}.
	\end{equation*}
	By substituting the currents within the inductor current balance equation, we obtain
	\begin{equation*}
		\frac{(U_\mathrm{1}-u_\mathrm{ds,on} - u_\mathrm{fw} - U_\mathrm{2}) \cdot T_\mathrm{on} }{L}=\frac{(U_\mathrm{2} + 2 u_\mathrm{fw}) \cdot T_\mathrm{off}}{L}.
	\end{equation*}
	Replacing $T_\mathrm{on}$ and $T_\mathrm{off}$ by the switching period yields to
	\begin{equation*}
		\frac{(U_\mathrm{1}-u_\mathrm{ds,on} - u_\mathrm{fw} - U_\mathrm{2}) \cdot T_\mathrm{s} \cdot D_\mathrm{1}}{L} =
		\frac{(U_\mathrm{2} + 2 u_\mathrm{fw}) \cdot (1-D_\mathrm{1}) \cdot T_\mathrm{s}}{L}.
	\end{equation*}
	After simplifying by canceling $T_\mathrm{s}$ and $L$, we obtain
	\begin{equation*}
		(U_\mathrm{1}-u_\mathrm{ds,on} - u_\mathrm{fw} - U_\mathrm{2}) \cdot D_\mathrm{1} =
		(U_\mathrm{2} + 2 u_\mathrm{fw}) \cdot (1-D_\mathrm{1}).
	\end{equation*}
	Solving this equation with respect to $D_\mathrm{1}$ yields
	\begin{equation*}
		D_\mathrm{1} = \frac{U_\mathrm{2}+ 2 u_\mathrm{fw}}{U_\mathrm{1}-u_\mathrm{ds,on} + u_\mathrm{fw}}.
	\end{equation*}
	For the operation points $U_\mathrm{1}=\SI{7.5}{\volt}$ and $U_\mathrm{1}=\SI{10}{\volt}$ we obtain
	\begin{align*}
		D_\mathrm{1, U_\mathrm{1} = \SI{7.5}{\volt}} & =\frac{\SI{5}{\volt} + 2 \cdot \SI{0.8}{\volt}}{\SI{7.5}{\volt} - \SI{0.5}{\volt} + \SI{0.8}{\volt}} = 0.85, \\
		D_\mathrm{1, U_\mathrm{1} = \SI{10}{\volt}}  & =\frac{\SI{5}{\volt} + 2 \cdot \SI{0.8}{\volt}} {\SI{10}{\volt} - \SI{0.5}{\volt} + \SI{0.8}{\volt}} = 0.64.
	\end{align*}

\end{solutionblock}

